\chapter{Applicazioni}
\noindent 
In questo capitolo vedremo alcune applicazioni pratiche dei concetti di statistica matematica studiati nei capitoli precedenti.

\section{Regioni di confidenza}
\noindent
Imaginemos que tenemos una moneda trucada donde sale cara con probabilidad $p\in ]0,1[$, ¿como puedo estimar $p$? Sabemos que 
$$ \lim_{n\to \infty} \dfrac{\#\text{caras}}{\#\text{lanzamientos}} = p \quad \text{c.s.}$$
y otra forma de hacerlo es lanzar la moneda $n$ veces y contar el número de caras obtenidas, entonces una estimación natural de $p$ es $p \approx \dfrac{\#\text{caras}}{\#\text{lanzamientos}}$. \\

\noindent
Sin embargo, esta afirmación tiene dos problemas:
\begin{itemize}
    \item no tiene en cuenta una posibilidad (improbable) que esté muy alejada de la media, como por ejemplo, que salgan todo caras.
    \item no cuantifica el error en función como de grande es $n$.
\end{itemize}

\begin{definicion}[Modello statistico]
    Un \tbf{modelo estadístico} es una terna $(X, \mathcal{F}, P_{\theta}: \theta \in \Theta)$, donde 
    \begin{itemize}
        \item $X$ es un conjunto.
        \item $\mathcal{F}$ es una $\sigma$-algebra sobre $X$.
        \item $(P_{\theta} : \theta \in \Theta)$ es una familia de medidas (al menos dos) de probabilidad sobre $(X, \mathcal{F})$ indexada por un conjunto de parámetros $\Theta$.
    \end{itemize}
\end{definicion}

\begin{ejemplo}
    En el ejemplo de la moneda trucada con $n$ lanzamientos tendríamos el modelo estadístico
    $$(X, \mathcal{F}, P_{p}: p \in ]0,1[)$$
    donde
    \begin{itemize}
        \item $X = \{0,\ldots, n\}$ es el conjunto de posibles números de caras obtenidas en $n$ lanzamientos.
        \item $\mathcal{F} = \mathcal{P}(X)$ es la $\sigma$-álgebra de partes de $X$.
        \item $P_{\theta}=B_{n,\theta}$, donde recordamos que 
        $$P(X=k)=B_{n,\theta}(\{k\})= \binom{n}{k} \theta^k (1-\theta)^{n-k}$$
        y tenemos $\Theta = ]0,1[$.
    \end{itemize}
\end{ejemplo}

\begin{definicion}[Regione di confidenza]
    Dado un modelo estadístico $(X, \mathcal{F}, P_{\theta}: \theta \in \Theta)$ y sea  $\tau: \Theta \to \Sigma$ una característica del sistema que se quiere estimar, donde $\Sigma$ es un conjunto arbitrario, entonces \tbf{una región de confianza para $\tau$ con nivel de error $\alpha\in ]0,1[$ (o nivel de confianza $1-\alpha$)} es una aplicación
    \begin{equation*}
        C: X \to \mathcal{P}(\Sigma) \quad \text{ tal que } \quad \inf_{\theta \in \Theta} P_{\theta}(\{x \in X : \tau(\theta) \in C(x)\}) \geq 1 - \alpha
    \end{equation*}
    Si $\Sigma=\R$ y $C(x)$ es un intervalo, se dice que $C$ es un \tbf{intervalo de confianza}.
\end{definicion}

\begin{observacion}
    Se puede usar indistintamente ``'\textit{nivel de error $\alpha$}'' o ``\textit{nivel de confianza $1-\alpha$}''.
\end{observacion}

\noindent
Un \underline{caso trivial} que podemos estudiar, es el caso en el que $C(x) = \Sigma$ para todo $x \in X$, ya que en este caso se cumple que
\begin{equation*}
    P_{\theta}(\{x \in X : \tau(\theta) \in C(x)\}) = P_{\theta}(\{x \in X : \tau(\theta) \in \Sigma\}) =P_{\theta}(\{x\in X\})= P_{\theta}(X) = 1 \geq 1 - \alpha \quad \forall \theta \in \Theta
\end{equation*}
por lo que en tal caso $\forall \alpha \in ]0,1[$, $C:X\to \{\Sigma\}$ es una región de confianza para $\tau$ con nivel de error $\alpha$. \\ 
\begin{observacion}
    Si $\tau$ no se especifíca explícitamente, se asume que es la función identidad, es decir, 
    $$\tau: \Theta \to \Sigma=\Theta \ \text{ con } \ \tau(\theta) = \theta \quad \forall \theta \in \Theta$$
\end{observacion}

\subsection{Metodo diretto per costruire regione di confidenza}
\begin{teo}[Costruzione di intervalli di confidenza]
    Dado $(X, \mathcal{F}, P_{\theta}: \theta \in \Theta)$ un modelo estadístico y $\tau=id: \Theta \to \Theta$, la identidad, una característica del sistema que se quiere estimar. Entonces se siguen estos pasos para construir un intervalo de confianza para $\tau$ con nivel de error $\alpha \in ]0,1[$:
    \begin{itemize}
        \item Para cada $\theta \in \Theta$ encontrar el $C_{\theta}$ más pequeño tal que $P_{\theta}(C_{\theta}) \geq 1 - \alpha$ y $C_{\theta} \subseteq X$.
        \item $C: X \to \mathcal{P}(\Theta)$ está dada por $C(x) = \{\theta \in \Theta : x \in C_{\theta}\}$.
    \end{itemize}
\end{teo}

\noindent
Podemos ahora preguntarnos si realmente esta construcción de la aplicación $C$ nos da un intervalo de confianza para $\tau$ con nivel de error $\alpha$, lo que significa ver si cumple que 
$$\inf_{\theta \in \Theta} P_{\theta}(\{x \in X : \tau(\theta) \in C(x)\}) \geq 1 - \alpha$$ 
Para ellos tomamos $\theta \in \Theta$ fijo, y tenemos que
\begin{align*}
    y\in \{x\in X : \tau(\theta)=\theta \in C(x)\} \iff  \theta \in C(y)=\{\theta\in \Theta \mid y\in C_{\theta}\} \iff y\in C_{\theta} 
\end{align*}
por lo que se cumple que $\{x\in X : \theta \in C(x)\} = C_{\theta}$, y por tanto para cada $\theta \in \Theta$ se cumple que
\begin{equation*}
    P_{\theta}(\{x \in X : \tau(\theta) \in C(x)\}) = P_{\theta}(C_{\theta}) \geq 1 - \alpha
\end{equation*}
es decir, se cumple $\inf\limits_{\theta \in \Theta} P_{\theta}(\{x \in X : \tau(\theta) \in C(x)\}) \geq 1 - \alpha$, de manera que la construcción es totalmente correcta, pues nos da $C$ que es un intervalo de confianza para $\tau$ con nivel de error $\alpha$.

\subsubsection{Intervalli di confidenza delle Binomiali}
\noindent
A continuación estudiaremos el caso particular de las distribuciones binomiales, y veremos cómo construir un intervalo de confianza para la característica del sistema $\tau=id: \Theta \to \Theta$, justificandolo mediante dos formas diferentes. \\
\noindent
Sea $(X, \mathcal{F}, P_{\mu}: \theta \in \Theta)$ un modelo estadístico donde 
\begin{itemize}
    \item $X=\{0,\ldots, n\}$ cuenta las veces que salió cara en $n$ lanzamientos de una moneda trucada.
    \item $P_{\theta}=B_{n,\theta}$ es la distribución binomial con parámetros $n$ y $\theta$.
    \item $\Theta=]0,1[$.
    \item $\tau=id: \Theta \to \Theta$ es la identidad, una característica del sistema que se quiere estimar.
\end{itemize}
entonces un intervalo de confianza para $\tau$ con nivel de error $\alpha \in ]0,1[$ está dado por la aplicación $C:X \to \mathcal{P}(\Theta)$ definida como
\begin{equation*}
    C(x)=\left(\dfrac{x}{n} -\veps, \dfrac{x}{n} + \veps\right)
\end{equation*}
de manera que se intenta aproximar $\theta$ por la frecuencia relativa $x/n$, ya que sabemos que la mejor aproximación de $\theta$ es la frecuencia relativa de caras obtenidas en los lanzamientos, por lo que se busca un intervalo centrado en dicho punto. Buscaremos ahora como calcular el valor de $\veps>0$, mediante dos métodos. \\ 

\noindent
\tbf{Chebyshev.} Sabiendo la restricción que debe cumplir $P_{\theta}$, tenemos lo siguiente
\begin{align*}
    P_{\theta}\left(\left\{x \in X : \theta \in C(x)\right\}\right) = P_{\theta}\left(\left\{x \in X : \theta \in \left(\dfrac{x}{n} -\veps, \dfrac{x}{n} + \veps\right)\right\}\right) = B_{n,\theta}\left(\left\{x \in X : \theta\in \left( \dfrac{x}{n} -\veps, \dfrac{x}{n} + \veps\right)\right\}\right) 
\end{align*}
\noindent
veámos a continuación cómo simplificar la expresión para $\theta$
\begin{align*}
    \theta \in \left(\dfrac{x}{n} -\veps, \dfrac{x}{n} + \veps\right) \iff \dfrac{x}{n} - \veps < \theta < \dfrac{x}{n} + \veps \iff -\veps < \theta - \dfrac{x}{n} < \veps \iff \left|\theta - \dfrac{x}{n}\right| < \veps
\end{align*}
tenemos por tanto 
\begin{align*}
    & P_{\theta}\left(\left\{x \in X : \theta \in C(x)\right\}\right) = B_{n,\theta}\left(\left\{x \in X : \left|\theta - \dfrac{x}{n}\right| < \veps\right\}\right) = 1-B_{n,\theta}\left(\left\{x \in X : \left|\theta - \dfrac{x}{n}\right| \geq \veps\right\}\right)\geq 1-\alpha \iff \\
    & \iff B_{n,\theta}\left(\left\{x \in X : \left|\theta - \dfrac{x}{n}\right| \geq \veps\right\}\right) \leq \alpha
\end{align*}
Definiremos ahora $A_1, \ldots, A_n$ eventos independientes con $P(A_i)=\theta$ de forma que la v.a. $X$ viene dada como
\begin{equation*}
    X = \sum\limits_{i=1}^{n} \mathbbm{1}_{A_i} \sim B_{n,\theta} 
\end{equation*}
entonces tenemos que
\begin{equation*}
    B_{n,\theta}\left(\left\{x \in X : \left|\theta - \dfrac{x}{n}\right| \geq \veps\right\}\right) = P\left(\left|\dfrac{X}{n} - \theta\right| \geq \veps\right) \overset{(*)}{\leq} \dfrac{Var\left[\dfrac{X}{n}\right]}{\veps^2}= \dfrac{Var[X]}{n^2 \veps^2} 
\end{equation*}
donde en $(*)$ hemos usado la desigualdad de Chebyshev, y sabemos que $\E[X/n] = \theta$. Desarrollando la varianza de $X$, obtenemos su valor simplificado
\begin{equation*}
    Var[X]=Var[\sum\limits_{i=1}^{n} \mathbbm{1}_{A_i}] \overset{(*)}{=}\sum\limits_{i=1}^{n} Var[\mathbbm{1}_{A_i}] = \sum\limits_{i=1}^{n} \theta(1-\theta) = n \theta(1-\theta)
\end{equation*}
donde en $(*)$ hemos usado la independencia de los eventos $A_i$, y además es fácil ver que \mbox{$Var[\mathbbm{1}_{A_i}] = \theta(1-\theta)$}. Finalmente, obtenemos la desigualdad esperada para $\veps$
\begin{align*}
    P_{\theta}\left(\left\{x \in X : \theta \in C(x)\right\}\right)\geq 1-\alpha \iff B_{n,\theta}\left(\left\{x \in X : |\theta - \dfrac{x}{n}| \geq \veps\right\}\right) \leq \alpha \iff \dfrac{n \theta(1-\theta)}{n^2 \veps^2}=\dfrac{\theta(1-\theta)}{n \veps^2} \leq \alpha 
\end{align*}
Como no quiero que $\veps$ dependa de $\theta$, usaré la cota superior de $\theta(1-\theta)$ en $]0,1[$, que es $1/4$, por lo que obtenemos la siguiente cota para $\veps$
\begin{equation*}
    \dfrac{1}{4 n \veps^2} \leq \alpha \iff \veps \geq \dfrac{1}{2 \sqrt{n \alpha}}
\end{equation*}
Es importante destacar que el haber hecho esta cota extra y la anterior al tener en cuenta la varianza, hace que el intervalo de confianza obtenido sea más grande de lo necesario, pero a cambio no depende de $\theta$. Notemos ahora lo siguiente sobre el valor de $\veps$ obtenido:
\begin{itemize}
    \item Si $\alpha << 1$, entonces $\veps$ es grande, y el intervalo de confianza $C(x)$ es \underline{grande}, lo cual es lógico ya que si queremos un nivel de error muy pequeño, el intervalo de confianza debe ser grande para asegurar que contiene el valor real de $\theta$.
    \item Si $n$ es grande, entonces $\veps$ es pequeño, y el intervalo de confianza $C(x)$ es \underline{pequeño} como queremos, ya que si hacemos muchos lanzamientos, la frecuencia relativa se acerca mucho a $\theta$, por lo que el intervalo de confianza debe ser pequeño.
\end{itemize}

\noindent
\tbf{Teorema del Límite Central.} Recordemos que el TCL en el caso de la bernoulli nos dice que
\begin{equation*}
    \dfrac{1}{\sqrt{n}} \sum\limits_{i=1}^{n} \dfrac{\mathbbm{1}_{A_i} - \theta}{\sqrt{\theta(1-\theta)}} \xrightarrow{d} N(0,1) \quad \text{cuando } n \to \infty
\end{equation*}
lo que implica que para $\Phi(x)=\dfrac{1}{\sqrt{2\pi}} \int_{-\infty}^{x} e^{-t^2/2} dt$ tenemos
\begin{equation*}
    F_{S_n^*}\to \Phi \quad \text{uniformemente en } \R \text{ cuando }n \to \infty \qquad \text{ donde } S_n^* = \dfrac{1}{\sqrt{n}} \sum\limits_{i=1}^{n} \dfrac{\mathbbm{1}_{A_i} - \theta}{\sqrt{\theta(1-\theta)}}
\end{equation*}
Antes de iniciar con los cálculos para obtener la aproximación de $\veps$, veámos dos desarrollos importantes que usaremos después
\begin{itemize}
    \item  $S_n^*= \dfrac{1}{\sqrt{n}} \sum\limits_{i=1}^{n} \dfrac{\mathbbm{1}_{A_i} - \theta}{\sqrt{\theta(1-\theta)}} = \dfrac{1}{\sqrt{n \theta(1-\theta)}} \left(\sum\limits_{i=1}^{n} \mathbbm{1}_{A_i} - n \theta\right) = \dfrac{X - n \theta}{\sqrt{n \theta(1-\theta)}}$ 
    \item $\left| \dfrac{x}{n} - \theta \right| < \veps \iff \sqrt{\dfrac{n}{\theta(1-\theta)}} \left|\dfrac{x - n \theta}{n}\right| < \veps \sqrt{\dfrac{n}{\theta(1-\theta)}}  \overset{(*)}{\iff} \dfrac{1}{\sqrt{n}}\left| \dfrac{x-n\theta}{\sqrt{\theta(1-\theta)}} \right| < \veps \sqrt{\dfrac{n}{\theta(1-\theta)}}$
\end{itemize}
donde en $(*)$ hemos usado que $\dfrac{\sqrt{n}}{|n|} = \dfrac{1}{\sqrt{n}}$, ya que también sabemos que $n\in \N$. Ahora sí, podemos empezar a calcular la cota para $\veps$
\begin{align*}
    & 1-\alpha \leq B_{n,\theta}\left(\left\{x \in X : \left|\dfrac{x}{n} - \theta\right| < \veps\right\}\right) = B_{n,\theta}\left(\left\{x \in X : \dfrac{1}{\sqrt{n}}\left| \dfrac{x-n\theta}{\sqrt{\theta(1-\theta)}} \right| < \veps \sqrt{\dfrac{n}{\theta(1-\theta)}}\right\}\right)= \\
    & = P\left(\dfrac{1}{\sqrt{n}}\left| \dfrac{X-n\theta}{\sqrt{\theta(1-\theta)}} \right| < \veps \sqrt{\dfrac{n}{\theta(1-\theta)}}\right) = P\left(|S_n^*| < \veps \sqrt{\dfrac{n}{\theta(1-\theta)}}\right) \overset{TCL}{\approx} \\
    & \overset{TCL}{\approx}  \Phi\left( \veps \sqrt{\dfrac{n}{\theta(1-\theta)}}\right) - \Phi\left(- \veps \sqrt{\dfrac{n}{\theta(1-\theta)}}\right) = 2 \Phi\left( \veps \sqrt{\dfrac{n}{\theta(1-\theta)}}\right) - 1
\end{align*}
Tenemos finalmente que
\begin{align*}
    & 1-\alpha \leq 2 \Phi\left( \veps \sqrt{\dfrac{n}{\theta(1-\theta)}}\right) - 1 \implies \Phi\left( \veps \sqrt{\dfrac{n}{\theta(1-\theta)}}\right) \geq \dfrac{2-\alpha}{2} \implies \\
    & \implies \veps \sqrt{\dfrac{n}{\theta(1-\theta)}} \geq \Phi^{-1}\left(\dfrac{2-\alpha}{2}\right) \implies \veps \geq \Phi^{-1}\left(\dfrac{2-\alpha}{2}\right) \sqrt{\dfrac{\theta(1-\theta)}{n}}
\end{align*}
Como bien sabemos si queremos un nivel de error $\alpha$ muy bajo, debemos por tanto tener un valor de $\veps$ grande, es decir, un intervalor $C$ grande. Veámos que esto se cumple en la expresión obtenida, suponiendo que $\alpha \to 0$
\begin{equation*}
    2\Phi\left(\veps \sqrt{\dfrac{n}{\theta(1-\theta)}}\right) - 1 \approx 1 \iff \Phi\left(\veps \sqrt{\dfrac{n}{\theta(1-\theta)}}\right) \approx 1 \iff \veps \to +\infty
\end{equation*}
donde claramente para el último resultado hemos usado la definición de $\Phi$, es decir, de función de distribución de la normal estándar.

\begin{observacion}
    Asumimos que $C$ es una aplicación tal que $\{x\in X\mid \tau(\theta) \in C(x)\}\in \mathcal{F}$.
\end{observacion}

\subsection{Metodo alternativo per costruire regione di confidenza}

\noindent
Fijado el modelo estadístico $(X, \mathcal{F}, P_{\theta}: \theta \in \Theta)$, supongamos que existe una función $T_{\theta}$ tal que $P_{\theta}\circ T_{\theta}^{-1}$ es constante, ¿como podemos usar esta función para construir $C$ la región de confianza para $\tau$ con nivel de error $\alpha$? \\
\noindent
La idea es poder definir $Q=P_{\theta}\circ T_{\theta}^{-1}$, de manera que dado $\alpha\in ]0,1[$ y $I_{\alpha}$ tal que $Q(I_{\alpha}) \geq 1-\alpha$, entonces $P_{\theta}(T_{\theta}^{-1}(I_{\alpha})) = Q(I_{\alpha}) \geq 1-\alpha$, por lo que demostrando \mbox{$\{\theta \in \Theta : T_{\theta}(x) \in I_{\alpha}\}=T_{\theta}^{-1}(I_{\alpha})$} ya habríamos terminado. Veámos primero varias definiciones y enunciados que nos llevarán a dicho resultado. \\

\begin{definicion}[Pivot per una caratteristica]
    Dado un modelo estadístico $(X, \mathcal{F}, P_{\theta}: \theta \in \Theta)$ y sea  $\tau: \Theta \to \Sigma$ una característica del sistema que se quiere estimar, entonces se dice que $(Q, T_{\mu}:\mu\in \Sigma)$ es un \tbf{pivote para $\tau$} si
    \begin{itemize}
        \item $\forall \mu \in \Sigma$, $T_{\mu}: X \to \Sigma'$ es una función medible, donde $(\Sigma', \mathcal{F}')$ es un espacio medible.
        \item $Q=P_{\theta} \circ T_{\tau(\theta)}^{-1} \quad \forall \theta \in \Theta$.
    \end{itemize}
\end{definicion}

\begin{observacion}
    Importante notar que $Q$ no depende de $\theta \in \Theta$, es decir, para todo $\theta \in \Theta$ sabemos que $Q$ es la misma, y además es una medida de probabilidad sobre $(\Sigma', \mathcal{F}')$.
\end{observacion}
\noindent
El pivote es una herramienta que ``estandariza'' el problema para que siempre trabajes con la misma distribución $Q$, facilitando enormemente el cálculo del intervalo.
\begin{notacion}
    $T_{\tau(\theta)}^{-1}(A) = \{x \in X : T_{\tau(\theta)}(x) \in A\}$ para todo $A \in \mathcal{F}'$.
\end{notacion}

\begin{lema}
    Dado $(X, \mathcal{F}, P_{\theta}: \theta \in \Theta)$ un modelo estadístico. Sea $(Q, T_{\mu}:\mu\in \Sigma)$ un pivote para la característica $\tau: \Theta \to \Sigma$, $\alpha \in ]0,1[$ y $I\subseteq \Sigma'$ tal que $Q(I) \geq 1-\alpha$, entonces la aplicación
    \begin{equation*}
        C: X \to \mathcal{P}(\Sigma), \quad C(x) = \{\mu \in \Sigma : x\in T_{\mu}^{-1}(I)\}
    \end{equation*}
    es una región de confianza para $\tau$ con nivel de error $\alpha$.
    \begin{proof}
        % //TODO: copiar de Jesus y Cristian
    \end{proof}
\end{lema}

\begin{ejemplo}[Prodotto Gaussiano]
    Consideremos el modelo estadístico dado por el producto Gaussiano $n-$veces
    \begin{equation*}
        (X, \mathcal{F}, P_{\mu}: \mu \in \R) = ( \R^n, \mathcal{B}(\R^n), N_{m,v}^{\otimes n} : m \in \R,v>0)
    \end{equation*}
    como viene dado dicho modelo queremos obtener una región de confianza $1-\alpha$ para poder estimar la media y la varianza, por lo que tenemos que los parámetros son $\Theta = \R \times \R^+$, donde $\theta=(m,v)$ y la característica del sistema que queremos estimar es $\tau: \Theta \to \Sigma=\Theta$, con $\tau(\theta) = \theta$. \\ 
    Buscamos por tanto, si existe, un pivote para $\tau$, es decir, un mapeo $T_{\theta}$ tal que $N_{m,v}^{\otimes n} \circ T_{m,v}^{-1}$ sea constante para $m,v$. Una buena elección para $T_{m,v}$ es la siguiente
    \begin{equation*}
        T_{m,v}:\R^n\to \R^n \quad \text{ dada por } \quad T_{m,v}(x)=\left(\dfrac{x_i-m}{\sqrt{v}}\right)_{i=1,\ldots,n}=\dfrac{x-m}{\sqrt{v}}
    \end{equation*}
    Veámos ahora que efectivamente $N_{m,v}^{\otimes n} \circ T_{m,v}^{-1}=N_{0,1}^{\otimes n}$ para todo $m\in \R$ y $v>0$, y por tanto $Q=N_{0,1}^{\otimes n}$ no depende de $\theta$. Basta demostrarlo para $A_1, \ldots, A_n \in \mathcal{B}(\R)$, ya que los \underline{rectángulos generan} \underline{la $\sigma$-álgebra de Borel en $\R^n$}, de modo que tenemos
    \begin{align*}
        & N_{m,v}^{\otimes n} \circ T_{m,v}^{-1}(A_1 \times \ldots \times A_n) = N_{m,v}^{\otimes n}\left(\left\{x \in \R^n : T_{m,v}(x) \in A_1 \times \ldots \times A_n\right\}\right) = \\
        & = N_{m,v}^{\otimes n}\left(\left\{x \in \R^n : \dfrac{x_i - m}{\sqrt{v}} \in A_i, i=1,\ldots,n\right\}\right) = N_{m,v}^{\otimes n}\left(\prod\limits_{i=1}^{n} (\sqrt{v} A_i + m)\right) = \\
        & = \prod\limits_{i=1}^{n} N_{m,v}(\sqrt{v} A_i + m) = \prod\limits_{i=1}^{n} N_{0,1}(A_i) = N_{0,1}^{\otimes n}(A_1 \times \ldots \times A_n)
    \end{align*}
    NO ESTA TERMINADO
    % //TODO: copiar de Jesus y Cristian caso 1 y caso 2
\end{ejemplo}

\begin{observacion}
    En la demostración se ha usado una idea muy importante sobre $\sigma-$álgebras de Borel en $\R$, que es que los intervalos abiertos generan la $\sigma-$álgebra de Borel en $\R$, de ahí que baste demostrar la igualdad en los rectángulos (los generadores) para que se cumpla en toda la $\sigma-$álgebra de Borel en $\R^n$.
\end{observacion}